% ============================================================================
%  Appendix C -- Reproducibility protocol
%
%  One thing left to fill at submission: the commit hash below.  Create the tag
%  first, then read the hash off it:
%
%      git tag -a submission -m "PX915 report submission"
%      git push origin submission
%      git rev-parse --short submission
%
%  The clone instructions check out the tag, so they work even if the hash line
%  is not refreshed.  Refresh the package versions too if the environment moves.
% ============================================================================

\section{Reproducibility protocol}
\label{app:protocol}

This appendix is written so that a reader who has never seen the code can
reproduce the wave speed quoted in Section \ref{sec:uq} and check it against a
reference value.

\subsection*{What is being reproduced}

The canonical result is the transition wave speed at the protocol point
%
\begin{equation*}
  A = 1.0, \qquad \delta = 0.001, \qquad \alpha = 0.10,
  \qquad N = 200, \qquad \Delta t = 10^{-4}, \qquad T = 400,
\end{equation*}
%
with expected value
%
\begin{equation*}
  v = 0.2206 \;\pm\; 0.0001\ \text{(numerical)} \;\pm\; 0.0219\ \text{(parametric)}.
\end{equation*}
%
The numerical figure is the convergence floor from Section
\ref{sec:verification}; the parametric figure is the spread induced by
$\pm10\%$ manufacturing tolerances on $\delta$, $A$ and $\alpha$, and is the
number a reader reproducing a \emph{physical} device should expect to meet. A
rerun of the code should match the reference to the numerical figure, not the
parametric one.

\subsection*{Code and environment}

\begin{description}
  \item[Repository]
    \url{https://github.com/Stephang217/1D_Diaghphragms_Fluid_Model}
  \item[Release]
    tag \texttt{submission}
  \item[Commit]
    \texttt{0000000} % TODO: refresh from `git rev-parse --short submission`
  \item[Language]
    Python 3.14.7
  \item[Packages]
    \texttt{numpy} 2.3.4, \texttt{pandas} 2.3.3, \texttt{matplotlib} 3.10.7,
    \texttt{scipy} 1.16.2, \texttt{jupyter}.
    Pinned in \texttt{requirements.txt}; install with
    \texttt{pip install -r requirements.txt}.
\end{description}

Reproducing the headline number needs \texttt{numpy} alone. The other packages
are needed only to read the cached ensembles and redraw the figures, and
\texttt{scipy} only by \texttt{diaphragm\_metafluid.ipynb}, which fits a front
width.

No compiled extensions, no GPU, and no random number generation in the canonical
run itself: the simulation is deterministic, so a correct rerun reproduces the
reference to floating-point round-off rather than to a tolerance. Randomness
enters only in the ensembles of Section \ref{sec:uq}, where every seed is
recorded in the CSV header along with the package versions that ensemble was
generated under.

\subsection*{The single canonical run}

\begin{verbatim}
    git clone https://github.com/Stephang217/1D_Diaghphragms_Fluid_Model
    cd 1D_Diaghphragms_Fluid_Model
    git checkout submission
    pip install -r requirements.txt

    python3 -c "
    import sys; sys.path.insert(0, 'src')
    from model import run_sim
    v, v_pred, ratio = run_sim(A=1.0, delta=0.001, alpha=0.10,
                               N=200, T=400, dt=1e-4)
    print(f'v_sim={v:.6f}  v_pred={v_pred:.6f}  ratio={ratio:.6f}')
    "
\end{verbatim}

\noindent Expected output, to be compared against
\texttt{results/protocol\_point.csv}:

\begin{verbatim}
    v_sim=0.220560  v_pred=0.220647  ratio=1.000393
\end{verbatim}

\noindent This is a single-core run taking approximately $50$~s on the
development machine (Apple Silicon laptop).

\subsection*{Reproducing the uncertainty}

The parametric band requires the ensemble rather than the single run:

\begin{verbatim}
    python3 scripts/uq_sweep.py          # 256 train + 32 test
    python3 scripts/disorder_sweep.py    # per-diaphragm disorder
\end{verbatim}

\noindent Each script prints a projected runtime from a single timed evaluation
before launching its pool, and writes a CSV into \texttt{results/} whose header
records the seed, the sample count, the lattice settings, the package versions
and the wall time. Loading them requires \texttt{comment='\#'}:
\texttt{pd.read\_csv(path, comment='\#')}. Opening
\texttt{uq\_reproducible\_result.ipynb} and running it top to bottom then
regenerates the tolerance band, the $\delta A/\alpha$ collapse and the headline
number from those CSVs.

On the development machine the tolerance ensemble takes roughly $53$~min for the
256 training samples and a further $6$~min for the 32 test samples, both across
$8$~workers; the disorder ensemble takes about $4$~min.

\subsection*{Results that ship pre-computed}

Several sweeps are too slow to recompute casually and are committed as cached
CSVs alongside the scripts that generate them, following the convention used
throughout the notebooks: a cached file is loaded if present, and setting
\texttt{FORCE\_RERUN = True} recomputes it.

\begin{description}
  \item[\texttt{results/pinning\_boundary.csv}]
    the $(\kappa,\Pi)$ pinning boundary of Section \ref{sec:pinning}; five
    bisections of seven runs each, roughly $70$~min. Regenerate with
    \texttt{python3 scripts/pinning\_boundary.py}.
  \item[\texttt{results/closed\_form\_validation.csv}]
    the 160-point grid behind Appendix \ref{app:tanh}, separating the
    energy-balance error from the $\tanh$ approximation; about $6$~min.
    Regenerate with \texttt{python3 scripts/closed\_form\_check.py}.
  \item[\texttt{results/alpha\_low\_windows.csv}]
    the low-damping time-window test of Section \ref{sec:validity}; the runs at
    the smallest $\alpha$ are long because $\tilde\tau = 2/\Omega$ is large and
    the lattice must be sized to the faster underdamped front.
\end{description}

\noindent Every other figure in the paper recomputes from scratch in at most a
few minutes.
